\documentclass{article}
\usepackage{booktabs}

\title{Transformer-Based Crop Disease Diagnosis}
\author{Example Author}
\date{}

\begin{document}
\maketitle

\section{Introduction}
Crop disease diagnosis from leaf images remains difficult when visually similar diseases appear under different lighting conditions. We study whether a vision transformer can improve classification quality over a conventional CNN baseline on this task.

\section{Method}
Our model replaces the CNN feature extractor with a transformer encoder and keeps the same classifier head and training recipe as the baseline.

\section{Results}
Table~\ref{tab:main-results} reports the comparison between the CNN baseline and our transformer model.

\begin{table}[h]
\centering
\caption{Classification accuracy on the held-out test set.}
\label{tab:main-results}
\begin{tabular}{lc}
\toprule
Model & Top-1 Accuracy \\
\midrule
CNN baseline & 89.4 \\
Transformer & 92.1 \\
\bottomrule
\end{tabular}
\end{table}

The transformer improves performance because global-context modeling reduces confusion between visually similar disease patterns while retaining discriminative lesion structure. The supporting evidence for this mechanism-based explanation is provided in the paragraph immediately below Table~2.

\section{Conclusion}
The transformer model provides a simple accuracy improvement on this dataset.

\end{document}
